\section{Erd\H{o}s Problem \#371: independent continuation}
\noindent Date: 2026-09-23. Research attempt; no claim of a new published result.
\subsection{1) TARGET AND SOURCE AUDIT}
Does $\{n:P^+(n)<P^+(n+1)\}$ have natural density $1/2$?

All supplied versions considered: \texttt{371.tex}.
The source cites logarithmic averaging and almost-all-scale results. I test whether those conclusions can bridge the remaining limit gap using only properties of counting functions. They cannot, as the following explicit model demonstrates.
\subsection{2) WORK}
Let $M_j=2^{2^j}$ and define a set $A$ by starting with the even positive integers and adjoining every integer in $[M_j,2M_j]$. Its logarithmic density is $1/2$: the added harmonic mass up to $X$ is at most $O(\#\{j:M_j\le X\})=O(\log\log X)=o(\log X)$.
However, at $X=2M_j$, the latest block adds $M_j/2+O(1)$ formerly absent odd integers. All earlier blocks have total size $o(M_j)$, since $M_{j-1}=\sqrt{M_j}$. Thus
\[
{|A\cap[1,2M_j]|\over2M_j}\longrightarrow\frac34.
\]
At $X=M_j-1$, earlier additions have density $o(1)$, giving a subsequential limit $1/2$.

This model even has density $1/2$ at almost all logarithmic scales in the quantitative sense that for every fixed $\epsilon>0$, the scales with deviation greater than $\epsilon$ have logarithmic density zero. Indeed additions up to $X$ are bounded by a constant times the largest $M_j\le X$. A deviation forces $X\le C_\epsilon M_j$, so the exceptional scales lie in the union of intervals $[M_j,C_\epsilon M_j]$, each of bounded logarithmic length. Their count is $O(\log\log X)$.
Therefore the generic counting-function structure, even together with those two averaging properties, cannot yield a full natural-density limit.
\subsection{3) ADVERSARIAL VERIFICATION}
The set $A$ is an abstract model and is not a largest-prime-factor comparison set. The countermodel refutes an implication between density notions, not the conjecture. Earlier blocks were controlled in ordinary size as well as logarithmic mass.
\subsection{4) FINAL}
\textbf{UNRESOLVED.}
An arithmetic argument excluding sparse exceptional scales for the actual comparison set is still required. No new density bound for prime factors is established.
\subsection{5) SOURCES CHECKED}
No external theorem is used in the new argument. The supplied source was read in full; the saved Tao status snapshot is dated 2026-09-22.
\subsection{6) COMPLETION ESTIMATE}
This is a conservative subjective measure of progress toward the original question, not a probability of correctness or a claim that the remaining work is routine.
\textbf{COMPLETION: 5\%}
